\documentclass{article}
\usepackage[preprint]{neurips_2026}

\usepackage[utf8]{inputenc}
\usepackage[T1]{fontenc}
\usepackage{hyperref}
\usepackage{url}
\usepackage{booktabs}
\usepackage{amsfonts}
\usepackage{amsmath}
\usepackage{microtype}
\usepackage{xcolor}

% This is an ES 678 course paper, not a conference submission.
\makeatletter
\renewcommand{\@noticestring}{ES 678 Course Project}
\makeatother

\title{Replace with a specific, informative project title}

\author{%
  Student Name \\
  Roll Number: 00000000 \\
  ES 678 \\
  Indian Institute of Technology Gandhinagar
}

\begin{document}
\maketitle

\begin{abstract}
% In 100--150 words: state the problem, data, method, evaluation, and main
% finding. Do not claim more than your evaluation supports.
Write a concise summary of your project here.
\end{abstract}

\section{Problem and Motivation}
State the question in concrete terms. Define the prediction target or pattern
of interest, what one data row represents, and when a prediction would be made.
Explain why the question matters and what claim your data can support.

\section{Data and Preparation}
Identify the data source, sample size, important variables, and access or
sharing limits. Report the data-quality checks relevant to your question.
Describe preprocessing and explain how every learned transformation was fit
using training data only.

\section{Methods}
Define a simple baseline and the methods you compare. Explain the objective,
assumptions, and reason for each choice. Include the mathematical formulation
that is needed to understand your method; define every symbol you use.

\section{Evaluation Design}
Describe the train/validation/test split or cross-validation scheme before
reporting results. Explain how model selection and hyperparameter tuning were
kept separate from final testing. Define each metric and why it fits the task.

\section{Results and Error Analysis}
Compare against the baseline with a compact table or readable figure. Interpret
the results in context. Inspect representative errors or failure cases. Report
uncertainty or subgroup results when they matter and can be supported.

\section{Graduate Depth Analysis}
Present at least one substantive analysis connected to this same project, such
as a derivation, sensitivity or ablation study, calibration or uncertainty
analysis, nontrivial baseline comparison, failure-mode audit, or deployment-
risk analysis. State what this extra analysis teaches beyond the main score.

\section{Limitations and Conclusion}
State the most important limitation and what evidence would be needed to make a
stronger claim. Do not interpret predictive association as proof of causation.
Conclude with the main answer to your project question.

\section{Reproducibility and AI-Use Disclosure}
Give code and data links or access instructions, split details, and enough
information to follow the main result. Briefly state which AI tools you used,
how you used them, and which outputs you independently checked. Never publish
private data, passwords, or API keys.

\begin{thebibliography}{9}
\bibitem[Author(2026)]{replace}
Replace this example with complete references for methods, data, and external
material used in the project.
\end{thebibliography}

\end{document}
